\chapter{Decisões Arquiteturais}

As decisões arquiteturais foram tomadas com o objetivo de equilibrar simplicidade inicial, extensibilidade futura e isolamento de responsabilidades. A seguir são discutidos os principais pontos de projeto e o racional por trás de cada escolha.

\section{Por que RabbitMQ?}

RabbitMQ foi escolhido por ser um broker maduro, amplamente documentado e adequado para cenários de desacoplamento entre produção e consumo de tarefas. Em um sistema de análise de imagens, a mensageria assíncrona evita que a API fique presa ao tempo de execução do processamento.

\section{Por que MassTransit?}

MassTransit simplifica a integração com RabbitMQ ao abstrair tópicos recorrentes de publicação, consumo, configuração e retries. Ele reduz boilerplate, melhora a composição da solução e torna a evolução da topologia de mensagens mais previsível \cite{masstransitdocs}.

\section{Por que callbacks?}

Callbacks foram adotados para permitir integração eficiente com o backend solicitante sem depender de polling constante. Isso melhora a fluidez da integração e reduz consultas desnecessárias ao microserviço.

\section{Por que Python para treinamento?}

Python é amplamente utilizado para pesquisa, preparação de dados e treinamento de redes neurais, com ecossistema consolidado para ciência de dados e deep learning. A escolha favorece experimentação e acelera o ciclo de construção do modelo \cite{pytorchdocs}.

\section{Por que C# para inferência?}

C# foi escolhido para a API e a orquestração do microserviço por sua forte integração com ASP.NET Core, tipagem estável, ecossistema maduro de web services e facilidade de composição com infraestrutura corporativa. Isso permite expor uma API robusta enquanto a camada de IA é incorporada de forma controlada.

\section{Por que Clean Architecture?}

Clean Architecture organiza o sistema em torno do domínio e das interfaces, reduzindo o impacto de mudanças em mensageria, storage e persistência. Essa abordagem é especialmente útil em soluções que passam por transição tecnológica, como o caso deste microserviço.

\section{Por que Repository?}

O Repository foi introduzido para encapsular o acesso à persistência e permitir uma troca futura de armazenamento sem alterar a aplicação. Mesmo na versão em memória, ele já preserva a fronteira entre regra de negócio e mecanismo de armazenamento.

\section{Por que ModelLoader?}

O ModelLoader foi criado para evitar recarga repetida do modelo e isolar a responsabilidade de vida útil do artefato em memória. Essa decisão será crucial quando o modelo real for integrado, especialmente se houver custo elevado de inicialização.

\begin{figure}[H]
    \centering
    \fbox{\parbox{0.85\textwidth}{\centering Espaço reservado para a tabela ou mapa de decisões arquiteturais complementares.}}
    \caption{Resumo visual das decisões arquiteturais}
    \label{fig:decisoes-arquiteturais}
\end{figure}
